% ============================================================================
%  A/00-map.tex — bản đồ phần A
%  Ngày tạo: 2026-09-06 | Sửa gần nhất: 2026-09-06 | Ôn gần nhất: chưa
% ============================================================================
\documentclass[../english.tex]{subfiles}
\begin{document}
\section*{Bản đồ phần A --- Nền tảng nghĩa}

\begin{tomtat}
Bốn chương chia theo \emph{câu hỏi chúng trả lời}: nghĩa nằm ở đâu (1); biết một từ là biết gì và cần bao nhiêu từ (2); ngữ pháp tối thiểu để tách câu dài (3); đoán nghĩa từ lạ qua cấu tạo (4). Chuỗi bắt buộc là 1 \(\rightarrow\) 2. Chương 3 và 4 là công cụ, đọc được bất cứ lúc nào và nên quay lại khi thực sự vướng.
\end{tomtat}

\subsection*{A. Phần này có gì}
\begin{itemize}
\item \textbf{\code{01-nghia-nam-o-dau}} --- nghĩa từ điển so với nghĩa trong ngữ cảnh; denotation và connotation; vì sao hai từ ``đồng nghĩa'' không thay thế được cho nhau; ba tầng nghĩa mà người đọc phải phục dựng.
\item \textbf{\code{02-biet-mot-tu}} --- chín mặt của việc biết một từ; từ vựng chủ động và bị động; họ từ; tần suất và ngưỡng độ phủ 95\% với 98\%; bao nhiêu từ là đủ để đọc tài liệu ngành.
\item \textbf{\code{03-ngu-phap-toi-thieu}} --- mệnh đề chính và mệnh đề phụ; động từ hữu hạn; bổ nghĩa dính vào đâu; thì và thể mang nghĩa gì trong văn khoa học; danh hoá; mạo từ và số nhiều cho người Việt.
\item \textbf{\code{04-tu-nguyen-va-cau-tao}} --- gốc Latin và Hy Lạp; tiền tố, hậu tố và cách chúng đổi từ loại; đoán nghĩa từ kỹ thuật; những chỗ từ nguyên đánh lừa.
\end{itemize}

\subsection*{B. Vì sao chia như vậy}
Cách chia thông thường của sách tiếng Anh là theo \emph{kỹ năng} --- từ vựng, ngữ pháp, đọc, viết --- và cách chia đó khiến người học coi từ vựng là một danh sách phải nhớ. Bốn chương ở đây chia theo một trục khác: từ \emph{nguồn của nghĩa} (chương 1) tới \emph{cách lưu nghĩa} (chương 2), rồi tới hai công cụ phục dựng nghĩa khi đọc (chương 3 và 4). Trục này phản ánh đúng thứ tự bạn cần chúng khi ngồi trước một trang tài liệu khó.

\subsection*{C. Phụ thuộc và thứ tự đọc}
\begin{figure}[H]\centering
\begin{tikzpicture}[node distance=0.62cm and 1.05cm,
  m/.style={draw=steel,fill=bluebg,rounded corners=2pt,inner sep=5pt,align=center,font=\small,text width=2.9cm},
  o/.style={draw=violetdark,fill=violetbg,rounded corners=2pt,inner sep=5pt,align=center,font=\small,text width=2.9cm},
  ar/.style={-{Latex[length=2.2mm]},draw=steel,thick},
  ard/.style={-{Latex[length=2.2mm]},draw=tiercol,thick,dashed},
  lb/.style={font=\scriptsize\itshape,color=reddark,align=center,text width=2.4cm}]
\node[m] (c1) {\textbf{1.} Nghĩa nằm\\ở đâu};
\node[m,right=of c1] (c2) {\textbf{2.} Biết một từ\\là biết gì};
\node[o,below=1.1cm of c1] (c3) {\textbf{3.} Ngữ pháp\\tối thiểu};
\node[o,below=1.1cm of c2] (c4) {\textbf{4.} Từ nguyên\\và cấu tạo};
\node[draw=rulecol,fill=codebg,rounded corners=2pt,inner sep=5pt,align=center,font=\small,
      text width=2.9cm,right=of c2] (b) {Phần B\\đọc sâu};
\draw[ar] (c1) -- node[lb,above=0pt] {cách lưu nghĩa} (c2);
\draw[ar] (c2) -- (b);
\draw[ard] (c3) -- (c1);
\draw[ard] (c4) -- (c2);
\draw[ard] (c3.east) -- (c4.west);
\end{tikzpicture}
\caption{Phụ thuộc giữa bốn chương. Chuỗi bắt buộc là 1 \(\rightarrow\) 2: phải hiểu nghĩa được quyết định thế nào thì cách ghi nhớ từ ở chương 2 mới có lý. Hai chương tô tím là công cụ, đọc lướt trước rồi quay lại khi vướng câu dài hoặc từ lạ.}
\label{fig:map-A-en}
\end{figure}

\subsection*{D. Cốt lõi hay tham khảo}
\textbf{Cốt lõi:} toàn bộ chương 1; phần chín mặt của việc biết một từ và các con số ngưỡng độ phủ ở chương 2; phần tách mệnh đề và phần danh hoá ở chương 3. \textbf{Tham khảo:} danh sách gốc từ đầy đủ ở chương 4 --- học theo nhu cầu, mỗi lần vài gốc gặp thật trong tài liệu của bạn, thay vì học thuộc bảng.
\end{document}
